\documentclass[border=0pt]{standalone}
% The shared preamble of every docs diagram (docs/diagrams/<name>.tex).
% dvisvgm's driver first, so patterns and opacity survive into the SVG.
\def\pgfsysdriver{pgfsys-dvisvgm.def}
\usepackage[siunitx,RPvoltages]{circuitikz}
\usepackage{xcolor}
% Latin Modern: outline fonts for every glyph a label uses, µ included.
\usepackage[T1]{fontenc}
\usepackage{lmodern}
\usetikzlibrary{arrows.meta,backgrounds,calc,patterns,positioning}
\tikzset{
  % A white card, so diagrams read in mdBook's dark themes too.
  every picture/.append style={show background rectangle,
    background rectangle/.style={fill=white,inner sep=4pt}},
  % Layout geometry: layers told apart by fill and hatching, never colour alone.
  layer a/.style={draw=black,line width=0.6pt,fill=gray!20},
  layer b/.style={draw=black,line width=0.6pt,fill=gray!45},
  layer c/.style={draw=black,line width=0.6pt,fill=white,postaction={pattern=north east lines}},
  layer d/.style={draw=black,line width=0.6pt,fill=white,postaction={pattern=crosshatch dots}},
  marker/.style={draw=black,dashed,line width=0.6pt},
  % A measured distance and its label.
  dim/.style={{Latex[length=2mm]}-{Latex[length=2mm]},line width=0.5pt},
  dim label/.style={font=\small,fill=white,inner sep=1pt},
  % Where a rule fails: thick, dashed and labelled, not only red.
  violation/.style={draw=red!75!black,line width=1.4pt,dash pattern=on 3pt off 2pt},
  violation label/.style={font=\small\bfseries,text=red!75!black},
  caption/.style={font=\small},
}

\begin{document}
\begin{circuitikz}[x=3cm,y=3cm]
  % The two-metal example layout, 1 unit = 1 um: met1 wires A (y 0..0.2)
  % and B (y 0.5..0.6), a via1 cut on A, and a met2 strap over both.
  \fill[layer a] (0,0) rectangle (2,0.2);
  \fill[layer a] (0,0.5) rectangle (2,0.6);
  \fill[layer c,opacity=0.85] (1.75,-0.2) rectangle (2,1);
  \fill[layer b] (1.8,0.025) rectangle (1.95,0.175);
  \draw (1.8,0.025) -- (1.95,0.175) (1.8,0.175) -- (1.95,0.025);
  % Labels.
  \fill (0.1,0.1) circle (0.6pt);
  \node[caption,anchor=west] at (0.12,0.1) {text \texttt{A} on met1\_label};
  \fill (0.1,0.55) circle (0.6pt);
  \node[caption,anchor=west] at (0.12,0.55) {text \texttt{B}};
  % Callouts.
  \node[caption,anchor=east] at (-0.05,0.1) {met1};
  \node[caption,anchor=east] at (-0.05,0.55) {met1};
  \node[caption,anchor=south] at (1.875,1) {met2};
  \draw[-{Latex[length=2mm]}] (2.3,0.1) node[caption,anchor=west] {via1 joins met1 to met2} -- (1.96,0.1);
  \draw[-{Latex[length=2mm]}] (2.3,0.55) node[caption,anchor=west,align=left] {met2 crosses B; no via,\\so B stays its own net} -- (2.01,0.55);
  \node[caption,anchor=north west,align=left] at (0,-0.3) {net A: met1 wire A, via1, met2 strap.\quad net B: met1 wire B.};
\end{circuitikz}
\end{document}
